\chapter{sklearn 接口与可复现流程速查}
\section{统一模型接口}
scikit-learn是本书的机器学习软件包，Python导入名为\texttt{sklearn}。其中\textbf{估计器}是支持\texttt{fit}等统一操作的程序对象；它可以是预测模型，也可以是从数据学习处理规则的变换器。\textbf{接口}指调用对象的方法与参数约定，不代表各对象内部执行同一算法。下表是学过正文后的查询材料，不要求在第一章记住所有类名。

sklearn 估计器通常通过 \texttt{fit(X, y)} 拟合，通过 \texttt{predict(X)} 预测。无监督方法可能只需 \texttt{fit(X)}，也可能提供 \texttt{fit\_predict(X)}。概率分类器常提供 \texttt{predict\_proba}，但不同模型并不一定具备这个接口。\texttt{score} 的默认定义随估计器类型变化，正式报告应直接调用明确命名的评价指标。
\section{预处理与 Pipeline}
标准化、填补、PCA 与模型组合时，使用 \texttt{Pipeline}。参数搜索把命名步骤与参数连接，例如 \texttt{clf\_\_C}。交叉验证中的每个训练折都应重新拟合整个流水线的可学习步骤。
\section{常见模型与任务对照}
\begin{longtable}{p{0.28\linewidth}p{0.35\linewidth}p{0.25\linewidth}}
\toprule
任务与方法&sklearn 类名&主要参数\\\midrule\endhead
线性回归&\texttt{LinearRegression}&\texttt{fit\_intercept}\\
Ridge/Lasso&\texttt{Ridge, Lasso}&\texttt{alpha}\\
逻辑回归&\texttt{LogisticRegression}&\texttt{C, max\_iter}\\
近邻分类&\texttt{KNeighborsClassifier}&\texttt{n\_neighbors}\\
决策树&\texttt{DecisionTreeClassifier}&\texttt{max\_depth}\\
随机森林&\texttt{RandomForestClassifier}&\texttt{n\_estimators}\\
梯度提升回归&\texttt{GradientBoostingRegressor}&\texttt{learning\_rate}\\
支持向量机&\texttt{SVC}&\texttt{C, gamma, kernel}\\
多层感知机&\texttt{MLPClassifier}&\texttt{hidden\_layer\_sizes}\\
降维&\texttt{PCA}&\texttt{n\_components}\\
聚类&\texttt{KMeans, DBSCAN}&\texttt{n\_clusters, eps}\\
异常检测&\texttt{IsolationForest}&\texttt{contamination}\\\bottomrule
\end{longtable}
\section{错误排查的推荐顺序}
出现错误时先检查数据与模型接口：样本数是否一致；训练与预测字段顺序是否一致；是否包含字符串、无穷或缺失值；是否在训练折拟合变换；参数名是否适用于当前版本。其次检查统计问题：类别是否严重不平衡；特征是否泄漏目标；测试集是否参与模型选择；评价指标是否与任务目标相符。
\section{独立验证的一个最小模板}
\begin{lstlisting}[caption={保留数据与开发集交叉验证的骨架}]
from sklearn.model_selection import train_test_split, cross_val_score
from sklearn.pipeline import make_pipeline
from sklearn.preprocessing import StandardScaler
from sklearn.linear_model import LogisticRegression

X_dev, X_test, y_dev, y_test = train_test_split(
    X, y, test_size=0.2, stratify=y, random_state=42)
pipe = make_pipeline(StandardScaler(), LogisticRegression(max_iter=2000))
print(cross_val_score(pipe, X_dev, y_dev, cv=5).mean())
# Holdout is used only after procedure and hyperparameters are fixed.
\end{lstlisting}


\section{三个不同层次的交叉验证}
外层最终保留用于最后检查，开发内部的交叉验证用于比较模型，内部早停或校准还可能再划分数据。每一层的“训练集”都是相对于该层而言。Pipeline能保证它所包裹的变换在外部交叉验证训练折中拟合，但不能自动理解模型内部再次划分的对象关系。

对分组数据，应把组标识传给需要它的分割器，并检查所有层都遵守任务要求。常见错误包括外层按对象分组、内层却随机拆行；或先用全部开发目标选择特征，再做看似严格的内部五折。发现这样的情况，应重新组织完整流程，而不是只在报告里补一句“采用交叉验证”。

\section{保存指标时保持方向一致}
scikit-learn的交叉验证通常约定分数越大越好，因此有些误差接口以负值返回，如 \nolinkurl{neg_root_mean_squared_error}。打印成RMSE时取一次负号即可；若随后又用最大值选择“最差误差”，容易选反。推荐把结果表列名明确写成 \texttt{rmse}、\texttt{accuracy}等，并在转换处检查一个小例子。

概率列顺序通过 \texttt{classes\_}核对；变换后的列顺序通过相关特征名接口核对。模型序列化时同时保存环境版本、特征清单和预处理，不从不可信来源加载任意序列化文件。不同版本间模型对象的可移植性不能只凭文件扩展名保证。

\section{随书实验的运行顺序}
在工程根目录安装列出的依赖后，运行 \texttt{python scripts/run\_deep\_experiments.py}执行17章深化实验及真实OGLE节选。该脚本按章运行并保存成功、警告和异常；一章失败不会被改名为通过。只重现真实OGLE节选可运行 \texttt{python experiments/project\_ogle.py}。所有这些输入都已包含或可以由明确规则生成。

重生成中文图形需要系统安装中文字体。源码会查找常见字体家族，工程不附带字体文件。已生成的PDF与PNG可直接用于LaTeX编译，即使本地暂未配置Python中文绘图字体也可以阅读和编译教材。

